\documentclass[12pt]{article}
\usepackage{amsfonts,amsthm,amsmath,amssymb,mathtools}
\usepackage{mathrsfs}
\usepackage{enumitem}
\usepackage{tikz-cd}
\usepackage{geometry}
\usepackage{mlmodern}

\geometry{letterpaper, portrait, margin=1in}

\setcounter{section}{0}

\renewcommand{\thesection}{\S\arabic{section}}
\renewcommand{\thesubsection}{\arabic{section}}
\newcommand{\dd}{\mathop{}\!\mathrm{d}}

\newtheorem{thm}{Theorem}
\newenvironment{theorem}[1]{
	\renewcommand{\thethm}{#1}
	\begin{thm}
		}{\end{thm}}

\newtheorem{lma}{Lemma}
\newenvironment{lemma}[1]{
	\renewcommand{\thelma}{#1}
	\begin{lma}
		}{\end{lma}}

\theoremstyle{definition}
\newtheorem{ex}{}
\newenvironment{exercise}[1]{
	\renewcommand{\theex}{#1}
	\begin{ex}
		}{\end{ex}}

\title{Homework 7}
\author{Connor Mooneyhan}
\date{October 16, 2025}

\begin{document}

\begin{center}
	{\bf\Large MTH 346/646 homework cover sheet}
\end{center}

\vspace{0.5in}

\noindent{\large Name}: Connor Mooneyhan \\
{\large Homework number}: 7 \\
{\large Date}: October 15, 2025 \\

{\Large What topics did this assignment cover? How comfortable do you feel with
each of those topics?}

Lutz-Nagell, it seems. I don't feel comfortable with the material at this point, so I'm ``taking the L'' on this homework, as they say, to give myself some time to try to go back and catch up since I feel like I've let myself slip through the cracks and don't really know what's going on at this point :(

Hoping next homework will go better!


\vspace{1.75in}

{\Large In the course of doing this assignment, where did you get stuck?
	How did you get unstuck?}

I attempted problem 1 and started writing some Magma code, but I realized I must be misunderstanding what $C(3^k)$ meant. I thought it meant perhaps that $nP$ was congruent to (0 : 1 : 0) mod $3^k$, but that must not be right.

\vspace{1.75in}

{\Large Is there anything you've done here that you're not confident about, that you want me to take a really close look at, or that you want me to suggest an alternate approach to?}

N/A

\pagebreak

\begin{ex}
	Let $C : y^{2} = x^{3} + 8x + 1$. Then $C(\mathbb{Q}) \cong \mathbb{Z}$ and
	a generator is $P = (2,5)$.

	For each $k$, $1 \leq k \leq 8$ determine the smallest positive integer $n$
	so that $nP \in C(3^{k})$.
\end{ex}
\begin{proof}
  I used this Magma code to attempt it\dots but I'm not sure I have the right idea of what $C(3^k)$ is since this gave no results for $n \leq 500$.
	\begin{verbatim}
    C := EllipticCurve([8, 1]);
    P := C![2, 5];

    for k in [1..8] do;
      printf "\n----------\nk=%o\n----------\n", k;
      modulus := 3^k;

      for i in [1..500] do;
        iP := i * P;
        // check if iP congruent to (0 : 1 : 0) mod 3^k
        if IsIntegral(iP[1]/modulus) and
           IsIntegral((iP[2] - 1)/modulus) and
           IsIntegral(iP[3]/modulus) then;
         printf "i=%o : %o\n", i, iP;
        end if;
      end for;
    end for;
  \end{verbatim}
\end{proof}

\begin{ex}
	Let $C : y^{2} = x^{3} + ax^{2} + bx + c$ and let $p$ be a prime
	for which $C/\mathbb{F}_{p}$ is non-singular mod $p$. Show that the reduction homomorphism
	$\phi_{p} : C(\mathbb{Q}) \to C(\mathbb{F}_{p})$ (defined in the day 9 notes) has kernel equal to
	$C(p)$.
\end{ex}

\begin{ex}
	$*$ Use the Lutz-Nagell theorem and problem 2 above to show that if $C : y^{2} = x^{3} + ax^{2} + bx + c$ is an elliptic curve with torsion subgroup isomorphic to $\mathbb{Z}/2\mathbb{Z} \times \mathbb{Z}/8\mathbb{Z}$,
	then $C/\mathbb{F}_{p}$ must be singular for $p = 2$, $p = 3$, $p = 5$, and $p = 7$.
\end{ex}

\end{document}
